% !PRINTED 277
\weberpage{277}
%<<CONT
und wollen nun nachweisen, dass, wenn wir die Voraussetzung hinzufügen, \gesperrt{dass von den Grössen (3) keine zwei auf einander folgende zugleich verschwinden}, wir eine \gesperrt{Sturm}'sche Kette vor uns haben.

Es ist dies eine einfache Folgerung aus der Formel, die wir schon im \S.\,82 zu einem ähnlichen Zwecke benutzt haben, und die wir so darstellen können:

\weq{4}{L_k S_k - T_k^2 = L_{k-1} L_{k+1},}

worin $S_k$ und $T_k$ ganze rationale Functionen von $x$ sind.

Wir haben nur zu zeigen, dass die Forderungen \S.\,86, 1. bis 4. befriedigt sind. Davon ist aber 1. in die Voraussetzung aufgenommen, 2. ist erfüllt, da das letzte Element gleich $\pm 1$ ist. Aus der Formel (4) folgt, dass, wenn $L_k=0$ ist, $L_{k-1}$ und $L_{k+1}$ entgegengesetzte Zeichen haben, dass also die Bedingung 3. erfüllt ist.

Es bleibt nur noch die Bedingung 4. übrig. Wenden wir aber die Formel (4) auf $k=n-1$ an, so lautet sie

\[
\frac{\partial L_n}{\partial a_{n,n}} \cdot \frac{\partial L_n}{\partial a_{n-1,n-1}} - \left(\frac{\partial L_n}{\partial a_{n,n-1}}\right)^2 = L_n L_{n-2},
\]

und daraus folgt, dass, wenn $L_n=0$ ist,

\[
\frac{\partial L_n}{\partial a_{n,n}}, \quad \frac{\partial L_n}{\partial a_{n-1,n-1}}
\]

gleiche Zeichen haben; und ebenso kann man schliessen, dass alle Haupt-Unterdeterminanten von $L_n$

\[
\frac{\partial L_n}{\partial a_{i,i}}
\]

dasselbe Zeichen haben.

Nun ist aber die Derivirte von $L_n$ (vgl. \S.\,22)

\[
L_n'(x) = -\sum_{1,n}^{i} \frac{\partial L_n}{\partial a_{i,i}},
\]

und hat also für einen Werth $x$, für den $L_n(x)$ verschwindet, das entgegengesetzte Zeichen, wie $L_{n-1}$, was eben die Forderung 4. verlangt. Daraus folgt der Satz:

I.\ \gesperrt{Sind $\alpha, \beta$ zwei reelle Werthe, $\alpha<\beta$, so ist die Anzahl der zwischen $\alpha$ und $\beta$ gelegenen Wurzeln von $L_n(x)$ gleich dem Ueberschuss der Anzahl der Zeichenwechsel der Kette (3) für $x=\alpha$ über die Anzahl der Zeichenwechsel für $x=\beta$.}

Die höchste Potenz von $x$, die in $L_k(x)$ vorkommt, ist, wie schon oben bemerkt,
%CONT>>
